% !TeX root = ../../../金融经济学_王江_讲义及习题.tex
% 教学补图：固定收益分布、r_F=0；CARA常投资额与CRRA常投资比例的示例。
\begin{center}
\begin{minipage}{\linewidth}
\centering
\begin{tikzpicture}[>=Stealth,font=\small,x=1cm,y=1cm]
 \draw[->] (0,0)--(4.3,0) node[right] {$w$};
 \draw[->] (0,0)--(0,3.1) node[above] {$a^*(w)$};
 \draw[blue!70!black,very thick] (.6,1)--(3.8,1);
 \draw[red!70!black,thick,dashed] (.6,.42)--(3.8,2.66);
 \node[blue!70!black,below] at (2.8,1) {CARA：常金额};
 \node[red!70!black,above left] at (3.8,2.66) {CRRA：同比增长};
 \node at (2,-.6) {（a）投资金额};
\end{tikzpicture}\hspace{1cm}
\begin{tikzpicture}[>=Stealth,font=\small,x=1cm,y=1cm]
 \draw[->] (0,0)--(4.3,0) node[right] {$w$};
 \draw[->] (0,0)--(0,3.1) node[above] {$a^*(w)/w$};
 \draw[blue!70!black,very thick,domain=.6:3.8,samples=60] plot (\x,{1/\x});
 \draw[red!70!black,thick,dashed] (.6,.7)--(3.8,.7);
 \node[blue!70!black,above right] at (.9,1.11) {CARA：比例下降};
 \node[red!70!black,above] at (2.7,.7) {CRRA：常比例};
 \node at (2,-.6) {（b）投资比例};
\end{tikzpicture}
\par\smallskip\small
固定收益分布、取 $r_F=0$，在各自适用的内点财富区间内，示例分别取 CARA 的 $a^*=1$ 和 CRRA 的 $a^*=0.7w$。CRRA 也具有 DARA：金额随财富上升，比例仍保持不变。
\WJFigure{fig:wj-investment-amount-share}{风险投资金额与投资比例}
\end{minipage}
\end{center}
